% Generated draft: requires mathematical and source review.
\documentclass[11pt]{article}
\usepackage[margin=25mm]{geometry}
\usepackage{fontspec}
\setmainfont{Latin Modern Roman}
\usepackage{amsmath,amssymb,mathtools,unicode-math}
\setmathfont{Latin Modern Math}
\usepackage{enumitem,needspace,xurl,hyperref}
\setlength{\parindent}{0pt}
\setlength{\parskip}{5pt}
\setlength{\emergencystretch}{3em}
\newcommand{\N}{\mathbb{N}}
\newcommand{\Z}{\mathbb{Z}}
\newcommand{\Q}{\mathbb{Q}}
\newcommand{\R}{\mathbb{R}}
\newcommand{\C}{\mathbb{C}}
\newcommand{\F}{\mathbb{F}}
\newcommand{\A}{\mathbb{A}}
\newcommand{\PP}{\mathbb{P}}
\newcommand{\E}{\mathbb{E}}
\newcommand{\eps}{\varepsilon}
\newcommand{\dd}{\,\mathrm{d}}
\begin{document}
\section*{UnsolvedMath ID: 30005435}
\textbf{Source problem code:} catalogue entry
\subsection{Definitions and mathematical statement}
A projective plane is an incidence structure \(\Pi=(P,\mathcal L,I)\) in which any two distinct points lie on exactly one common line, any two distinct lines meet in exactly one point, and there exist four points no three of which are collinear. Its subspace lattice \(L(\Pi)\) consists of the empty subspace, the points, the lines, and the whole plane, ordered by incidence and inclusion. The meet is intersection and the join is the smallest subspace containing both arguments.

For a division ring \(D\), let \(\operatorname{PG}(2,D)\) be the incidence plane of one-dimensional and two-dimensional right \(D\)-linear subspaces of \(D^3\), with incidence given by inclusion. A projective plane is Desarguesian if it is isomorphic to \(\operatorname{PG}(2,D)\) for some division ring \(D\), and non-Desarguesian otherwise.

The supplied question asks whether there exist a non-Desarguesian projective plane \(\Pi\) and a brick \(B\) whose underlying lattice is isomorphic to \(L(\Pi)\). Schematically, writing \(U(B)\) for that underlying lattice, the question is whether
\[
\exists\,\Pi\;\exists\,B\quad
\bigl(\Pi\text{ is a non-Desarguesian projective plane}\bigr)
\ \land\ 
\bigl(B\text{ is a brick}\bigr)
\ \land\ 
U(B)\cong L(\Pi).
\]
Here the isomorphism is a lattice isomorphism. The supplied record does not define a brick or its underlying-lattice construction. Consequently, this display records the intended existence question but is not a complete axiomatic formulation: the predicate defining bricks and the construction \(U\) must be recovered from the source before mathematical review. No restriction to finite planes is stated in the supplied question; whether an ambient finiteness convention applies remains unresolved.

The original extraction continues with adjacent material on \(\wedge\)-semibraces. It defines a \(\wedge\)-semibrace to be a set \(X\) with binary operations \(\cdot\) and \(\wedge\), and a distinguished element \(e\in X\), such that \(\wedge\) is associative, commutative, and idempotent, and, for all \(x,y,z\in X\),
\[
\begin{aligned}
x\cdot(y\wedge z)&=(x\cdot y)\wedge(x\cdot z),\\
e\cdot x&=x,\\
(x\wedge y)\cdot z&=(x\cdot y)\cdot(x\cdot z),\\
x\cdot e&=e.
\end{aligned}
\]
The excerpt attributes this definition to its reference [6]. Its subsequent Problem 3 asks which semilattices underlie a \(\wedge\)-semibrace, and which underlie one that is also an \(L\)-algebra. These are adjacent questions, not supplied definitions of a brick; no equivalence with the plane-lattice question is asserted here.
\subsection{Short English statement}
Does there exist a non-Desarguesian projective plane whose lattice of subspaces is the underlying lattice of a brick? The supplied excerpt does not define a brick, so a complete mathematical formulation requires additional source context. The record labels the problem open; that status is supplied curation, not an independently verified conclusion.
\subsection{Sources}
\begin{itemize}
\item Local UnsolvedMath registry: supplied complete record for permanent ID \(30005435\), titled \emph{Non-Desarguesian Plane Lattices as Bricks}, in the Oberwolfach Reports Open Problems collection. Exact source problem code: \texttt{catalogue entry}. The supplied proposed year is 2023; no individual proposer is identified.
\item \emph{Mini-Workshop: Skew Braces and the Yang--Baxter Equation} (2023), Oberwolfach Reports, pp.~537--563, DOI \texttt{10.4171/owr/2023/9}. Supplied source URL and literature-source link, labeled ``original OWR report'': \url{https://doi.org/10.4171/owr/2023/9}. This report citation is supplied provenance; it was not independently consulted for this draft. The adjacent semibrace definition cites reference [6], whose bibliographic details are absent from the supplied record.
\item Supplied literature assessment, checked \texttt{2026-08-22}: the report presents the non-Desarguesian-plane brick question as open and records only positive constructions for Desarguesian subspace lattices; the supplied assessment reports that no primary source located in that curation supplies a brick on a non-Desarguesian plane lattice. This is supplied dated curation, not independently verified. The background explicitly calls it dated literature triage and requests independent verification of time-sensitive status and cited claims.
\item The supplied registry marks the statement \texttt{corrected\_verified} and reports that the official report was checked directly and that the clean statement preserves the mathematical target while removing extraction artifacts. These are supplied verification claims, not verification performed for this draft. In particular, the supplied excerpt still lacks a definition of a brick.
\end{itemize}
\end{document}
